
%%%%%%%%%%%%%%%%%%%%%%%%%%%%%%%%%%%%%%%%%%%%%%%%%%%%%%%%%%%%%%%%%%%%%%%%%%%%%%%%
% littlefs: rbyd B-trees: efficient B-trees when RAM << block size
%%%%%%%%%%%%%%%%%%%%%%%%%%%%%%%%%%%%%%%%%%%%%%%%%%%%%%%%%%%%%%%%%%%%%%%%%%%%%%%%

\documentclass[letterpaper,twocolumn,10pt]{article}
\usepackage{usenix-2020-09}

% newlines in tables
\usepackage{makecell}
\renewcommand\theadfont{\normalsize}

% remove padding around lists
\usepackage{enumitem}
\renewcommand{\descriptionlabel}[1]{\hspace\labelsep\normalfont\bfseries #1.}
\setlist[description]{
    nosep,
    wide,
%    labelindent=0.5\parindent,
    labelindent=0pt,
%    leftmargin=0pt,
%    leftmargin=0pt,
%%    itemindent=2em,
%    itemsep=0pt,
%    topsep=1\parsep,
%    parsep=1\parsep,
%    partopsep=0pt,
    %parsep=0pt plus 4pt minus 0pt,
    %itemsep=0pt plus 4pt minus 0pt,
}
\setlist[enumerate]{
    label={\arabic*.},
    font=\normalfont\bfseries,
%    label={\arabic*.},
    nosep,
%    wide,
    labelindent=0pt,
%    leftmargin=1em,
%    itemindent=2em,
%%    leftmargin=1.5em,
%    itemsep=0pt,
%    topsep=1\parsep,
%    parsep=1\parsep,
%    partopsep=0pt,
    %parsep=0pt plus 4pt minus 0pt,
    %itemsep=0pt plus 4pt minus 0pt,
}



% new list!
% 
% no idea how to make \item implicit :/
%
% also no idea why \denumeratei breaks :/
%
% it's all an ugly mess :/
%
\newcounter{ditemi}
\renewcommand{\theditemi}{(\number\value{ditemi})}
\newcommand{\denumeratereset}[0]{
    \setcounter{ditemi}{0}
}
\newcommand{\denumerateitem}[1]{
    \normalfont\bfseries(\number\value{ditemi})~#1.
}
\newcommand{\ditem}[1][]{
    \refstepcounter{ditemi}
    \item[#1]
}
\newlist{denumerate}{enumerate}{1} % last arg is max-depth
\setlist[denumerate]{
    nosep,
    labelindent=0pt,
    %labelindent=-0.5\parindent,
    align=left,
    itemindent=!,
    before=\denumeratereset,
    format=\denumerateitem,
}

% prevent footnote splitting, write smaller footnotes!
% \interfootnotelinepenalty=10000


% try to condense things
% TODO we may need to remove these
\usepackage{microtype}
\usepackage[compact]{titlesec}
% \setlength{\intextsep}{3pt} % this is for in-text floats? ohhh
% \setlength{\textfloatsep}{3pt}
% \setlength{\dbltextfloatsep}{3pt}
% \setlength{\abovecaptionskip}{2pt plus 2pt minus 1pt}
% \setlength{\belowcaptionskip}{2pt plus 2pt minus 1pt}
\setlength{\textfloatsep}{3pt plus 12pt minus 0pt}
\setlength{\dbltextfloatsep}{3pt plus 12pt minus 0pt}
\setlength{\abovecaptionskip}{3pt}
\setlength{\belowcaptionskip}{3pt}
% \setlength{\abovedisplayskip}{0pt} % these are for maths?
% \setlength{\belowdisplayskip}{0pt}
% \setlength{\abovedisplayshortskip}{0pt}
% \setlength{\belowdisplayshortskip}{0pt}

% fit more figures
% \setcounter{topnumber}{4}
% \setcounter{totalnumber}{6}
% \renewcommand{\topfraction}{.9}
% \renewcommand{\textfraction}{.1}


% to be able to draw some self-contained figs
\usepackage{tikz}
% TODO correct?
\usetikzlibrary{
    positioning,
    calc,
    arrows,
    arrows.meta,
    shapes.geometric,
    patterns,
    patterns.meta,
    decorations.pathreplacing,
    automata,
    pgfplots.groupplots}
\usepackage{pgfplots}
%\usepackage{pgffor}
\usepackage{xstring}
\usepackage{subcaption} 
\usepackage{amsmath}
\usepackage{amssymb}
%\usepackage{mathtools}
\usepackage{graphicx}
\usepackage{ninecolors}
\usepackage{ifthen}
\usepackage{booktabs}
\usepackage{multirow}
\usepackage{refcount}
\usepackage{verbatim}
\usepackage{hyphenat}
\usepackage{xurl}

%% cache tikz figures
%\usetikzlibrary{external}
%\tikzexternalize[prefix=build/]

% pgfcompat
\pgfplotsset{compat=1.18}

% allow hyphenating texttt
%\DeclareFontFamily{\encodingdefault}{\ttdefault}{\hyphenchar\font=`\-}
\DeclareFontFamily{\encodingdefault}{\ttdefault}{\hyphenchar\font=`\-}

% hyphenation rules
% \hyphenation{micro-con-troller}


% some conveniences
\newcommand{\reffig}[1]{Figure~\ref{fig:#1}}
\newcommand{\refeval}[1]{Figure~\ref{eval:#1}}
\newcommand{\reftbl}[1]{Table~\ref{tbl:#1}}
\newcommand{\refsec}[1]{\S\ref{sec:#1}}
\newcommand{\refeqn}[1]{Equation~\eqref{eqn:#1}}

% subfigure counter
\newcounter{subfig}[figure]
\renewcommand{\thesubfig}{\the\numexpr\value{figure}+1(\alph{subfig})}
\newcommand{\sublabel}[2]{
    \refstepcounter{subfig}
    \llap{(\alph{subfig})}
    \label{#1:\the\numexpr\value{figure}+1:#2}
}
\newcommand{\subfig}[1]{\sublabel{fig}{#1}}
\newcommand{\subeval}[1]{\sublabel{eval}{#1}}
\newcommand{\refsub}[3]{\hyperref[#1:#2]{\ref*{#1:\getrefnumber{#1:#2}:#3}}}
\newcommand{\refsubfig}[2]{Figure~\refsub{fig}{#1}{#2}}
\newcommand{\refsubeval}[2]{Figure~\refsub{eval}{#1}{#2}}

% hack from https://tex.stackexchange.com/a/698806
%
% I don't actually know how this works!
%
\makeatletter
\def\footlabel{%
    \edef\@currentHref{footnote.\the\c@footnote}%
    % adjust aim up
    \raisebox{\baselineskip}[0pt]{\hypertarget{\@currentHref}{}}%
    \label%
}
\makeatother

% old habits die hard
\def\lt{<}
\def\gt{>}

% some colors
\colorlet{rbydred}{red!20}
\colorlet{rbydgreen}{green!20}
\colorlet{rbydblue}{blue!20}
\colorlet{rbydyellow}{yellow!30}
\colorlet{rbydmagenta}{magenta!30}
\colorlet{rbydcyan}{cyan!30}

\colorlet{rbyddarkred}{red!60!black}
\colorlet{rbyddarkgreen}{green!60!black}
\colorlet{rbyddarkblue}{blue!60!black}
\colorlet{rbyddarkyellow}{yellow!50!black}
\colorlet{rbyddarkmagenta}{magenta!50!black}
\colorlet{rbyddarkcyan}{cyan!50!black}

\colorlet{rbydlightred}{red!12}
\colorlet{rbydlightgreen}{green!12}
\colorlet{rbydlightblue}{blue!12}
\colorlet{rbydlightyellow}{yellow!12}
\colorlet{rbydlightmagenta}{magenta!12}
\colorlet{rbydlightcyan}{cyan!12}



%%% YES yes controls %%%

% anonymize for blind review!
% \def\yesanon{}

% comment these out to speed things up
%
\def\yesinputfig{}
\def\yesinputeval{}

% unfortunately not everything can make the cut...
% \def\yeschop{}


% some aliases
\def\lfs+{littlefs}
% anonymize!
\ifdefined\yesanon
\def\lfsiii+{avianfs}
\else
\def\lfsiii+{littlefs3}
\fi
\def\lfsii+{littlefs2}
\def\spiffs+{SPIFFS}
\def\yaffsii+{Yaffs2}

% anything with a prefix also needs to be anonymized
\ifdefined\yesanon
\def\lfsiiifilefruncate+{\texttt{avfs\_file\_fruncate}}
\def\lfsiiifsgc+{\texttt{avfs\_fs\_gc}}
\else
\def\lfsiiifilefruncate+{\texttt{lfs3\_file\_fruncate}}
\def\lfsiiifsgc+{\texttt{lfs3\_fs\_gc}}
\fi

% some other shortcuts
\def\lfsiiimetadatamax+{\texttt{metadata\_max}}
\def\lfsiiiblockrecycles+{\texttt{block\_recycles}}
\def\lfsiiishrubsize+{\texttt{shrub\_size}}
\def\lfsiiicachesize+{\texttt{cache\_size}}
\def\lfsiiifragmentsize+{\texttt{fragment\_size}}
\def\lfsiiicrystalthresh+{\texttt{crystal\_thresh}}
\def\lfsiiilookaheadsize+{\texttt{lookahead\_size}}
\def\lfsiiigbmaprebuildthresh+{\texttt{gbmap\_rebuild\_thresh}}
\def\lfsiiisplitthresh+{\texttt{split\_thresh}}

% Winbond NOR/NAND flash shortcuts
\def\wnor+{\texttt{W25Q64JV}}
\def\wnand+{\texttt{W25N01GV}}


% disable inputs to speed things up
\ifdefined\yesinputfig
    \newcommand{\inputfig}[1]{\input{#1}}
\else
    \newcommand{\inputfig}[1]{\texttt{#1.tex}}
\fi
\ifdefined\yesinputeval
    \newcommand{\inputeval}[1]{\input{#1}}
\else
    \newcommand{\inputeval}[1]{\texttt{#1.tex}}
\fi

% this one's always on and makes grepping easy
\newcommand{\inputwip}[1]{\input{#1}}

% chopping controls
\ifdefined\yeschop
    \newenvironment{chop}{}{}
\else
    \newenvironment{chop}{\comment}{\endcomment}
\fi


% gimme a line
\def\linewip{\par\noindent\rule{\linewidth}{0.5pt}}




%%%%%%%%%%%%%%%%%%%%%%%%%%%%%%%%%%%%%%%%%%%%%%%%%%%%%%%%%%%%%%%%%%%%%%%%%%%%%%%%
\begin{document}
%%%%%%%%%%%%%%%%%%%%%%%%%%%%%%%%%%%%%%%%%%%%%%%%%%%%%%%%%%%%%%%%%%%%%%%%%%%%%%%%

% don't want date printed
\date{}

% make title bold and 14 pt font (Latex default is non-bold, 16 pt)
% \title{\Large \bf
%     \lfs+: rbyd B-trees: Efficient B-trees when RAM $\ll$ block size
% }
% \title{\Large \bf
%     \lfs+: rbyd B-trees: Incremental B-trees when RAM $\ll$ block size
% }
\title{\Large \bf
    %\lfs+: rbyd B-trees: Flexible incremental B-trees when RAM $\ll$ block size
    %\lfs+: rbyd B-trees: Incremental B-trees when RAM $\ll$ block size
    \lfs+: rbyd B-trees: Flexible B-trees when RAM $\ll$ block size
}
% \title{\Large \bf
%     \lfs+: rbyd B-trees: Incremental order-statistic B-trees when RAM $\ll$
%     block size
% }

\ifdefined\yesanon\else
%for single author (just remove % characters)
\author{
{\rm Christopher Haster}\\
\lfs+
% copy the following lines to add more authors
% \and
% {\rm Name}\\
%Name Institution
} % end author
\fi

\maketitle



%-------------------------------------------------------------------------------
\section{Introduction}
\label{sec:introduction}
%-------------------------------------------------------------------------------

\lfs+ is a little filesystem targeting microcontrollers---small devices with
%{$\sim$}256KiB--1MiB of ROM, {$\sim$}32KiB--128KiB of RAM, but---thanks to
MiBs of ROM, KiBs of RAM, and, thanks to
increasing storage density, potentially GiBs of storage.
%with proportionally large block sizes.
With block sizes in the range of {$\sim$}128KiB--1MiB,
block-based CoW algorithms are slow, and
RAM-dependent algorithms impossible.
% block sizes in the range of {$\sim$}128KiB--1MiB.
% These block sizes make block-based CoW algorithms slow, and
% RAM-dependent algorithms impossible.
% %To make matters worse,
% If this wasn't enough of a challenge,
% increasing storage density is also reflected in
% increasing block sizes.
% %NAND flash's {$\sim$}128KiB--1MiB blocks, for example, make
% %block-based CoW algorithms slow, and RAM-dependent algorithms impossible.
% On the upper end, the {$\sim$}128KiB--1MiB blocks found in NAND flash
% make block-based CoW algorithms slow, and RAM-dependent algorithms
% impossible.
% %With NAND flash block sizes in the
% %{$\sim$}128KiB--1MiB range, block-based CoW algorithms are slow, and
% %RAM-dependent algorithms impossible.
% %Given these constraints,
To support increasingly dense storage, we are trying to find a flexible
B-tree with the following constraints:
% To support increasingly dense storage, we are trying to find an
% order-statistic~\cite{TODO} B-tree with the following constraints:
\begin{enumerate}
\item Strictly bounded-RAM, i.e. no non-tail recursion
\item RAM usage independent of block size
\item Does not rely on full block rewrites
\end{enumerate}
% \item Support for sparse keys, e.g. file names
% \item Support for order-statistic keys, e.g. compressable ranges
% \end{enumerate}
These constraints rule out traditional array-backed
B-trees~\cite{ubiquitous-btree}, %which rewrite full blocks,
which rely on rewriting full blocks,
% during updates,
and
% rely on rewriting full blocks during updates, and
%the log-backed B-trees found in bcachefs~\cite{bcachefs}, which rely on
log-backed B-trees such as those found in bcachefs~\cite{bcachefs},
which rely on
%bcachefs's log-backed B-trees~\cite{bcachefs}, which rely on
reconstructing logs in-RAM.
% log-backed
% B-trees, such as those found in bcachefs~\cite{bcachefs}, which rely on
% reconstructing logs in-RAM.
% which support
% incremental updates, but rely on reconstructing logs in RAM.

% These constraints rule out existing B-tree designs.
% Traditional B-trees~\cite{ubiquitous-btree}
% implement inner nodes using dense arrays, and while
% this may be optimal in terms of branching-factor and storage usage,
% updates require rewriting full blocks.
% Log-backed B-trees, such as those found in bcachefs~\cite{bcachefs},
% support incremental updates, but rely on reconstructing logs in RAM.

Previous versions of \lfs+ avoided B-trees for these reasons~\cite{littlefs2},
however this has proven problematic for performance~\cite{li-r-27}.
% Fortunately, the Dhara FTL~\cite{dhara}
% %---a lightweight FTL also targeting microcontrollers---
% has shown it's possible to encode
In \lfsiii+, we are exploring an alternative B-tree
%that leverages
built on
log-encoded binary trees pioneered by the Dhara FTL~\cite{dhara}.
This so-called
\emph{red-black-yellow Dhara (rbyd) B-tree}
% \emph{rbyd B-tree}
provides a
flexible incremental B-tree in bounded-RAM,
with support for both sparse keys and order-statistic ranges.
% with support for both sparse
% keys, e.g. file names, and order-statistic keys~\cite{TODO, TODO}, e.g.
% compressable ranges.
% This so-called \emph{red-black-yellow Dhara (rbyd) B-tree} provides an
% incremental order-statistic B-tree in bounded-RAM.
%referred to as a \emph{red-black-yellow Dhara (rbyd) B-tree}. Rbyd B-trees 
%provide a flexible, incremental B-tree in bounded-RAM.


% % \begin{enumerate}
% % \item Strictly bounded-RAM, i.e. no non-tail recursion
% % \item RAM usage independent of block size
% %% \item RAM usage does not scale with block size
% % \item Incremental, i.e. does not rely on full block rewrites
% % \end{enumerate}
% 
% 
% 
% %This creates unique challenges for filesystem design.
% %(1)~Libraries are required to run in bounded RAM, limiting
% %data-structures to strictly tail-recursive algorithms.
% %(2)~Storage block sizes frequently exceed the
% %device's available RAM, making the reading of a full block into RAM
% %impossible.
% %This creates a unique set of requirements for filesystem design:
% %\begin{enumerate}
% %\item Strictly bounded RAM, i.e. no non-tail recursion
% %\item Limited caching
% %\item 
% %\end{enumerate}
% This creates a unique challenge for
% filesystem design: How do you design a filesystem to work efficiently in
% bounded-RAM?
% To make matters worse, increasing storage density is reflected in
% increasing block sizes.
% With block sizes in the {$\sim$}128KiB--1MiB range, NAND flash, for
% example, makes block-based CoW algorithms slow, and RAM-dependent
% algorithms impossible.
% Given these constraints, how 
% % creates a
% % performance nightmare for block-based CoW algorithms, and an
% % impossibility for algorithms that rely on reading blocks into RAM.
% % NAND flash, for example, often uses a block size of
% % {$\sim$}128KiB--1MiB. impossible to fit in some device's memory, and a
% % % TODO nightmare? bottleneck?
% % performance nightmare for block-based CoW algorithms.
% 
% % provides a {$\sim$}128KiB--1MiB block
% % size.
% % %currently the densest form of microcontroller storage,
% % implements a block size in the {$\sim$}128KiB--1MiB range.
% % 
% % 
% % With block sizes in the {$\sim$}128KiB--1MiB range for NAND flash, 
% 
% 
% %(1)~strictly
% %bounded-RAM, i.e. no non-tail recursion, and (2)~blocks that may not fit in
% %a devices working memory.
% % with {$\sim$}4KiB--128KiB--1MiB blocks.
% %This creates a unique challenge for
% %filesystem design: How do you design a filesystem to work efficiently in
% %bounded-RAM?
% 
% %\begin{description}
% %
% %\item[Bounded RAM] Libraries are required to run in bounded RAM, limiting
% %data-structures to strictly tail-recursive algorithms.
% %
% %\item[RAM {$\ll$} block size] Storage block sizes frequently exceed the
% %device's available RAM, making the reading of a full block into RAM
% %impossible.
% %
% %\end{description}
% 
% % An important piece of this puzzle is finding flexible data-structures that
% % can be reused to minimize code cost. For this purpose, B-trees are extremely
% % enticing. Unfortunately, existing designs rely on reading a full block into
% % RAM, which is simply not possible for many of these devices.
% % \lfs+ currently uses {$\sim$}2KiB of RAM, so this is an unreasonable cost even
% % for NOR flash's relatively small 4KiB blocks.
% % 
% % Previous versions of \lfs+ avoided B-trees for this reason, however this
% % proved problematic for performance.
% % % limited \lfs+'s performance.
% % In \lfsiii+ we are exploring \emph{rbyd B-trees},
% % a modification of the traditional B-tree algorithm in which inner nodes are
% % implemented using red-black-yellow Dhara trees (rbyds), self-balancing
% % log-encoding trees based on the Dhara FTL~\cite{dhara}.
% % Relying entirely on tail-recursive algorithms, rbyd B-trees promise to provide
% % a flexible B-tree in strictly bounded RAM.
% 
% %In previous versions of \lfs+, we avoided B-trees for this reason, instead
% %relying on simple skip-lists and linked-lists. However, these data-structures
% %are limited and pushed us to explore B-trees more thoroughly.
% %
% %In \lfsiii+ we are exploring the adoption of an \emph{rbyd B-trees}, the
% %composition of an in-block red-black-yellow Dhara tree (rbyd) and B-tree, that
% %provides an efficient B-tree without needing to load full blocks into RAM.
% 
% % TODO wait, this is not necessarily the case
% Traditional B-trees use densely packed arrays for their inner
% nodes---dense arrays are optimal in terms of branching factor and storage
% usage, but require full block rewrites to update.
% bcachefs~\cite{TODO} adopted a variant where inner nodes are backed by
% logs. This allows for incremental B-tree updates, but requires
% reconstructing logs in-RAM.
% 
% 
% requires reading a full block into RAM.
% Previous versions of \lfs+~\cite{TODO,TODO} avoided B-trees for these reasons,
% however this has proven problematic for performance.
% In \lfsiii+, we are revisiting B-trees, this time exploring a B-tree variant
% that leverages log-encoded trees pioneered by the Dhara FTL~\cite{dhara}---%
% a lightweight FTL also targeting microcontrollers.


%-------------------------------------------------------------------------------
\section{Design}
\label{sec:design}
%-------------------------------------------------------------------------------

The core idea is to replace the inner nodes of traditional
B-trees~\cite{ubiquitous-btree} with rbyd trees---%
%\emph{red-black-yellow Dhara (rbyd)} trees---%
a self-balancing order-statistic variant of the Dhara tree implemented in the
Dhara FTL~\cite{dhara}.
The result is a flexible B-tree, with inner nodes that can be
manipulated incrementally without being read into RAM.

The original Dhara tree implements a radix tree, and while it is
possible to extend the radix tree into a B-tree, the resulting
B-tree is relatively inflexible.
Augmenting a traditional B-tree with sparse keys and order-statistic ranges
does not change the $O(b) \rightarrow O(b)$ runtime for block size $b$,
but attempting the same transformation on the emulated arrays in a Dhara
B-tree increases the
runtime $O(\log b) \rightarrow O(b \log b)$.

To solve this, we need to extend the sparse and order-statistic
properties into the Dhara tree itself.
Our solution is a \emph{red-black-yellow Dhara (rbyd) tree}, %\footnotemark,
in which we use colors to emulate the inner branches of a 2-3-4 tree,
mimicking a traditional red-black tree~\cite{rbtree}.
% Like a traditional red-black tree, the rbyd algorithm is tail-recursive
% and runs in bounded-RAM.
This provides self-balancing and order-statistic operations, while
remaining strictly tail-recursive.
% This implements a self-balancing order-statistic Dhara tree, and is
% strictly tail-recursive.

% \footnotetext{
%     We also explored red-black Dhara (rbd) B-trees, which support sparse keys,
%     but attempting to add order-statistic operations resulted in
%     ambiguities.
%     %in an ambiguous
%     %tree.
%     %We also explored a red-black Dhara (rbd) tree,
%     %which supports sparse
%     %keys, but an additional color was needed to support order-statistic
%     %operations.
% }

The extension of rbyds into a full B-tree is relatively straightforward.
Rbyds contain either leaf data or indirect B-tree pointers, and most
updates require only a cheap log append. When a log is full, we \emph{compact}
the relevant rbyd into a new block, at which point the CoW
semantics of B-trees take over.
%To avoid compacting too frequently, we split when a node is half-full, but
%only when the
%% relevant
%log needs compaction.
To avoid runaway compactions, we split when a node is half-full, but
only when compaction is necessary.
% but only check for this condition when the log is full-full.
% To avoid runaway compactions, we split when half-full, but we only check
% merge/split conditions during compactions.
This
%We wait until compaction to check
%merge/split conditions, which
%amortizes the relevant costs,
%amortizes merge/split checks,
amortizes compaction work,
introduces hysteresis, and results in
%better
optimal
%leaf distribution during
leaf distribution after
%sequential updates.
sequential writes.

% We also implement tail-recursive tree rebalancing during compactions,
% which brings the compacted storage cost down
% $O(b \log b) \rightarrow O(b)$.
% In practice, the overhead of rbyds divides the branching-factor by
% {$\sim$}8
% %when compared to dense arrays.
% over dense arrays.
% %of dense
% %arrays by {$\sim$}8.
% % divides the branching-factor by
% % {$\sim$}8.

The main downside of rbyds is increased storage overhead, however a
tail-recursive rebalancing algorithm brings compacted cost down
$O(b \log b) \rightarrow O(b)$.
% thanks to tail-recursive rebalancing algorithm, this is only a
% constant.
% brings the compacted overhead down
% to a constant, $O(b \log b) \rightarrow O(b)$.
% Rbyds have a higher storage overhead than dense arrays. However,
% tail-recursive rebalancing during compactions brings this down
% $O(b \log b) \rightarrow O(b)$.
In our current implementation, the overhead of rbyds divides the
%branching-factor by {$\sim$}8 over dense arrays.
branching-factor by ${\sim}1/8$ over dense arrays.

% Each leaf in a given rbyd points to either a child B-tree node
% Instead of checking for merge/split conditions every commit, 


% A traditional B-tree can be extended to support sparse keys and
% order-statistic ranges, but only because inner node operations already
% scale $O(b)$ for block size $b$. When emulating arrays, Dhara trees scale
% $O(\log b)$, but attempting to apply the same transformation pushes the total
% runtime up to $O(b\log b)$.
%
%If we attempted to extend Dhara trees
%the same way, the result would scale 
%
%Applying the same transformation to
%inner nodes backed by Dhara trees pushes the $O(\log b)$ update cost up to
%$O(b\log b)$.
%
%
%You may think that because traditional B-trees are implemented with
%arrays, a Dhara B-tree can be extended to support sparse keys and
%order-statistic ranges, but this is not true.

% ---
% 
% The original Dhara tree implements a radix tree, and, while a radix
% B-tree may be useful for many applications, it is insufficient for
% littlefs's use case.
% 
% We've actually found three possible B-tree variants, each with support
% for increasingly complex keys.
% 
% 
% % The result is a flexible B-tree structure with inner nodes that can be
% % manipulated without being fully read into RAM.
% 
% % The main challenge with adopting Dhara trees is the need for self-balancing
% % operations.
% 
% Unfortunately, we can't use the original Dhara tree as is, as it only
% implements a radix tree. This works well for Dhara FTL's dense
% virtual{$\rightarrow$}physical block mapping, but results in an unbalanced
% tree when keys are sparse.
% 
% 
% translation, but 
% 
% The original Dhara
% 
% In the Dhara FTL, the original Dhara tree is implemented as a radix tree. This
% works well for the Dhara FTL's virtual{$\rightarrow$}physical block mapping,
% but is limiting as a general purpose data-structure.



%The central idea is to replace the inner nodes of B-trees with a data-structure
%that can be manipulated without being fully read into RAM.
%A promising starting point in this area is the Dhara FTL~\cite{dhara}---%
%a lightweight FTL also targeting microcontrollers---which showed it is possible
%to tail-recursive update a tree encoded in a log by writing only $O(\log n)$
%pointers.

%to encode a radix tree in a log, and tail-recursive update the tree by writing
%$O(\log n)$ pointers.
%However, to make this useful for 

%Our starting point is a Dhara tree, the core data-structure of the
%Dhara FTL~\cite{dhara}, a lightweight FTL also targeting microcontrollers.
%However, the Dhara FTL only provides a radix tree, and 
%
%The Dhara FTL~\cite{dhara}---A lightweight FTL also targeting
%microcontrollers---showed it is possible to implement a log-encoded radix tree
%in bounded RAM. We will refer to this as a \emph{Dhara tree}.
%
%To make this useful for a multipurpose B-tree

% %-----------------------------------
% \subsection{Rebalancing compactions}
% %-----------------------------------
% 
% %-----------------------------------------------
% \subsection{Red-black-yellow Dhara (rbyd) trees}
% %-----------------------------------------------

%-------------------------------------------------------------------------------
\section{Results and Discussion}
\label{sec:results}
\label{sec:discussion}
%-------------------------------------------------------------------------------

Rbyd B-trees are currently implemented as
% primary
the main file data-structure in
\lfsiii+. They have survived significant testing %, including
with assertions
over the rbyd color-invariants, and simulated measurements show the
expected logarithmic runtimes.

Measuring simulated
% file
write throughput shows improved performance
characteristics around larger block sizes in \lfsiii+ when compared to
\lfsii+.
While \lfsii+'s performance degrades as block size increases, \lfsiii+'s
performance
% actually improves due to the increased branching-factor.
generally improves due to reduced metadata overhead.

However, the more complex file data-structure is not without downsides.
Static analysis shows \lfsiii+'s code increasing from
16.9KiB {$\rightarrow$} 34.5KiB, and RAM increasing from
2.0KiB {$\rightarrow$} 2.8KiB. % note this includes ctx+stack!
This includes unrelated changes, but we believe the added file
complexity is the main culprit.
% This includes unrelated changes in the filesystem, but we believe the
% added file complexity and advanced file operations are major culprits.

For future work, we will be looking to better understand overall
filesystem performance---focusing on throughput and latency, comparing 
against \lfsii+ and other filesystems in this space---%
as well as a deeper analysis into the increased code cost and possible
savings.


% more complex file data-structure is the main culprit.

% When measuring simulated file write throughput, \lfsiii+ shows much
% better performance around larger block sizes---as block size increases,
% \lfsii+'s write throughput degrades, while \lfsiii+'s actually increases
% due to the increased branching-factor.
% 
% 
% Measuring simulated file write throughput in both \lfsiii+ and \lfsii+
% shows 
% 
% Comparing simulated file write throughput in both \lfsiii+ and \lfsii+
% shows improved 
% We've measured simulated throughputs of sequential, random, and logging
% file operations for both \lfsiii+ and \lfsii+, which 
% 
% 
% 
% Static analysis shows \lfsiii+ growing 
% 
% Comparing 
% 
% In total, static analysis shows \lfsiii+ growing 
% 
% 
% , which we plan to stabilize this year.
% 
% data-structure,
% 
% So far, we have focused on static and simulated results. We compare against
% \lfsii+, the previous version of \lfs+, which uses a simpler, append-only
% CTZ skip-list~\cite{ctz-skip-list} data-structure.
% 
% Thanks to avoiding recursion, we are able to calculate an upper bound on stack
% usage, which shows an increase from 2.0KiB {$\rightarrow$} 2.9KiB.
% 
% % using models based on
% %popular user devices. 
% 
% Simulated throughput

% %-------------------------------------------------------------------------------
% \section{Discussion}
% \label{sec:discussion}
% %-------------------------------------------------------------------------------







%-------------------------------------------------------------------------------
\bibliographystyle{plain}
\bibliography{littlefs-btree-abstract}

%%%%%%%%%%%%%%%%%%%%%%%%%%%%%%%%%%%%%%%%%%%%%%%%%%%%%%%%%%%%%%%%%%%%%%%%%%%%%%%%
\end{document}
%%%%%%%%%%%%%%%%%%%%%%%%%%%%%%%%%%%%%%%%%%%%%%%%%%%%%%%%%%%%%%%%%%%%%%%%%%%%%%%%

